%!TeX root=../emmatop.tex
\chapter[Chapter \thechapter]{}
\lettrine[lraise=0.3]{E}{very} body in and about Highbury who had ever visited Mr Elton, was disposed to pay him attention on his marriage. Dinner-parties and evening-parties were made for him and his lady; and invitations flowed in so fast that she had soon the pleasure of apprehending they were never to have a disengaged day.

<I see how it is,> said she. <I see what a life I am to lead among you. Upon my word we shall be absolutely dissipated. We really seem quite the fashion. If this is living in the country, it is nothing very formidable. From Monday next to Saturday, I assure you we have not a disengaged day!—A woman with fewer resources than I have, need not have been at a loss.>

No invitation came amiss to her. Her Bath habits made evening-parties perfectly natural to her, and Maple Grove had given her a taste for dinners. She was a little shocked at the want of two drawing rooms, at the poor attempt at rout-cakes, and there being no ice in the Highbury card-parties. Mrs Bates, Mrs Perry, Mrs Goddard and others, were a good deal behind-hand in knowledge of the world, but she would soon shew them how every thing ought to be arranged. In the course of the spring she must return their civilities by one very superior party—in which her card-tables should be set out with their separate candles and unbroken packs in the true style—and more waiters engaged for the evening than their own establishment could furnish, to carry round the refreshments at exactly the proper hour, and in the proper order.

Emma, in the meanwhile, could not be satisfied without a dinner at Hartfield for the Eltons. They must not do less than others, or she should be exposed to odious suspicions, and imagined capable of pitiful resentment. A dinner there must be. After Emma had talked about it for ten minutes, Mr Woodhouse felt no unwillingness, and only made the usual stipulation of not sitting at the bottom of the table himself, with the usual regular difficulty of deciding who should do it for him.

The persons to be invited, required little thought. Besides the Eltons, it must be the Westons and Mr Knightley; so far it was all of course—and it was hardly less inevitable that poor little Harriet must be asked to make the eighth:—but this invitation was not given with equal satisfaction, and on many accounts Emma was particularly pleased by Harriet's begging to be allowed to decline it. <She would rather not be in his company more than she could help. She was not yet quite able to see him and his charming happy wife together, without feeling uncomfortable. If Miss Woodhouse would not be displeased, she would rather stay at home.> It was precisely what Emma would have wished, had she deemed it possible enough for wishing. She was delighted with the fortitude of her little friend—for fortitude she knew it was in her to give up being in company and stay at home; and she could now invite the very person whom she really wanted to make the eighth, Jane Fairfax.— Since her last conversation with Mrs Weston and Mr Knightley, she was more conscience-stricken about Jane Fairfax than she had often been.—Mr Knightley's words dwelt with her. He had said that Jane Fairfax received attentions from Mrs Elton which nobody else paid her.

<This is very true,> said she, <at least as far as relates to me, which was all that was meant—and it is very shameful.—Of the same age—and always knowing her—I ought to have been more her friend.—She will never like me now. I have neglected her too long. But I will shew her greater attention than I have done.>

Every invitation was successful. They were all disengaged and all happy.—The preparatory interest of this dinner, however, was not yet over. A circumstance rather unlucky occurred. The two eldest little Knightleys were engaged to pay their grandpapa and aunt a visit of some weeks in the spring, and their papa now proposed bringing them, and staying one whole day at Hartfield—which one day would be the very day of this party.—His professional engagements did not allow of his being put off, but both father and daughter were disturbed by its happening so. Mr Woodhouse considered eight persons at dinner together as the utmost that his nerves could bear—and here would be a ninth—and Emma apprehended that it would be a ninth very much out of humour at not being able to come even to Hartfield for forty-eight hours without falling in with a dinner-party.

She comforted her father better than she could comfort herself, by representing that though he certainly would make them nine, yet he always said so little, that the increase of noise would be very immaterial. She thought it in reality a sad exchange for herself, to have him with his grave looks and reluctant conversation opposed to her instead of his brother.

The event was more favourable to Mr Woodhouse than to Emma. John Knightley came; but Mr Weston was unexpectedly summoned to town and must be absent on the very day. He might be able to join them in the evening, but certainly not to dinner. Mr Woodhouse was quite at ease; and the seeing him so, with the arrival of the little boys and the philosophic composure of her brother on hearing his fate, removed the chief of even Emma's vexation.

The day came, the party were punctually assembled, and Mr John Knightley seemed early to devote himself to the business of being agreeable. Instead of drawing his brother off to a window while they waited for dinner, he was talking to Miss Fairfax. Mrs Elton, as elegant as lace and pearls could make her, he looked at in silence—wanting only to observe enough for Isabella's information—but Miss Fairfax was an old acquaintance and a quiet girl, and he could talk to her. He had met her before breakfast as he was returning from a walk with his little boys, when it had been just beginning to rain. It was natural to have some civil hopes on the subject, and he said,

<I hope you did not venture far, Miss Fairfax, this morning, or I am sure you must have been wet.—We scarcely got home in time. I hope you turned directly.>

<I went only to the post-office,> said she, <and reached home before the rain was much. It is my daily errand. I always fetch the letters when I am here. It saves trouble, and is a something to get me out. A walk before breakfast does me good.>

<Not a walk in the rain, I should imagine.>

<No, but it did not absolutely rain when I set out.>

Mr John Knightley smiled, and replied,

<That is to say, you chose to have your walk, for you were not six yards from your own door when I had the pleasure of meeting you; and Henry and John had seen more drops than they could count long before. The post-office has a great charm at one period of our lives. When you have lived to my age, you will begin to think letters are never worth going through the rain for.>

There was a little blush, and then this answer,

<I must not hope to be ever situated as you are, in the midst of every dearest connexion, and therefore I cannot expect that simply growing older should make me indifferent about letters.>

<Indifferent! Oh! no—I never conceived you could become indifferent. Letters are no matter of indifference; they are generally a very positive curse.>

<You are speaking of letters of business; mine are letters of friendship.>

<I have often thought them the worst of the two,> replied he coolly. <Business, you know, may bring money, but friendship hardly ever does.>

<Ah! you are not serious now. I know Mr John Knightley too well—I am very sure he understands the value of friendship as well as any body. I can easily believe that letters are very little to you, much less than to me, but it is not your being ten years older than myself which makes the difference, it is not age, but situation. You have every body dearest to you always at hand, I, probably, never shall again; and therefore till I have outlived all my affections, a post-office, I think, must always have power to draw me out, in worse weather than to-day.>

<When I talked of your being altered by time, by the progress of years,> said John Knightley, <I meant to imply the change of situation which time usually brings. I consider one as including the other. Time will generally lessen the interest of every attachment not within the daily circle—but that is not the change I had in view for you. As an old friend, you will allow me to hope, Miss Fairfax, that ten years hence you may have as many concentrated objects as I have.>

It was kindly said, and very far from giving offence. A pleasant <thank you> seemed meant to laugh it off, but a blush, a quivering lip, a tear in the eye, shewed that it was felt beyond a laugh. Her attention was now claimed by Mr Woodhouse, who being, according to his custom on such occasions, making the circle of his guests, and paying his particular compliments to the ladies, was ending with her—and with all his mildest urbanity, said,

\begin{figure}[tbph]
\centering
\includegraphics[width=.9\linewidth]{34rain}
\caption[I am very sorry to hear of your being out in the rain]{<I am very sorry to hear, Miss Fairfax, of your being out this morning in the rain.>}
\end{figure}

<I am very sorry to hear, Miss Fairfax, of your being out this morning in the rain. Young ladies should take care of themselves.—Young ladies are delicate plants. They should take care of their health and their complexion. My dear, did you change your stockings?>

<Yes, sir, I did indeed; and I am very much obliged by your kind solicitude about me.>

<My dear Miss Fairfax, young ladies are very sure to be cared for.—I hope your good grand-mama and aunt are well. They are some of my very old friends. I wish my health allowed me to be a better neighbour. You do us a great deal of honour to-day, I am sure. My daughter and I are both highly sensible of your goodness, and have the greatest satisfaction in seeing you at Hartfield.>

The kind-hearted, polite old man might then sit down and feel that he had done his duty, and made every fair lady welcome and easy.

By this time, the walk in the rain had reached Mrs Elton, and her remonstrances now opened upon Jane.

<My dear Jane, what is this I hear?—Going to the post-office in the rain!—This must not be, I assure you.—You sad girl, how could you do such a thing?—It is a sign I was not there to take care of you.>

Jane very patiently assured her that she had not caught any cold.

<Oh! do not tell me. You really are a very sad girl, and do not know how to take care of yourself.—To the post-office indeed! Mrs Weston, did you ever hear the like? You and I must positively exert our authority.>

<My advice,> said Mrs Weston kindly and persuasively, <I certainly do feel tempted to give. Miss Fairfax, you must not run such risks.—Liable as you have been to severe colds, indeed you ought to be particularly careful, especially at this time of year. The spring I always think requires more than common care. Better wait an hour or two, or even half a day for your letters, than run the risk of bringing on your cough again. Now do not you feel that you had? Yes, I am sure you are much too reasonable. You look as if you would not do such a thing again.>

<Oh! she shall not do such a thing again,> eagerly rejoined Mrs Elton. <We will not allow her to do such a thing again:>—and nodding significantly—<there must be some arrangement made, there must indeed. I shall speak to Mr E. The man who fetches our letters every morning (one of our men, I forget his name) shall inquire for yours too and bring them to you. That will obviate all difficulties you know; and from us I really think, my dear Jane, you can have no scruple to accept such an accommodation.>

<You are extremely kind,> said Jane; <but I cannot give up my early walk. I am advised to be out of doors as much as I can, I must walk somewhere, and the post-office is an object; and upon my word, I have scarcely ever had a bad morning before.>

<My dear Jane, say no more about it. The thing is determined, that is (laughing affectedly) as far as I can presume to determine any thing without the concurrence of my lord and master. You know, Mrs Weston, you and I must be cautious how we express ourselves. But I do flatter myself, my dear Jane, that my influence is not entirely worn out. If I meet with no insuperable difficulties therefore, consider that point as settled.>

<Excuse me,> said Jane earnestly, <I cannot by any means consent to such an arrangement, so needlessly troublesome to your servant. If the errand were not a pleasure to me, it could be done, as it always is when I am not here, by my grandmama's.>

<Oh! my dear; but so much as Patty has to do!—And it is a kindness to employ our men.>

Jane looked as if she did not mean to be conquered; but instead of answering, she began speaking again to Mr John Knightley.

<The post-office is a wonderful establishment!> said she.—<The regularity and despatch of it! If one thinks of all that it has to do, and all that it does so well, it is really astonishing!>

<It is certainly very well regulated.>

<So seldom that any negligence or blunder appears! So seldom that a letter, among the thousands that are constantly passing about the kingdom, is even carried wrong—and not one in a million, I suppose, actually lost! And when one considers the variety of hands, and of bad hands too, that are to be deciphered, it increases the wonder.>

<The clerks grow expert from habit.—They must begin with some quickness of sight and hand, and exercise improves them. If you want any farther explanation,> continued he, smiling, <they are paid for it. That is the key to a great deal of capacity. The public pays and must be served well.>

The varieties of handwriting were farther talked of, and the usual observations made.

<I have heard it asserted,> said John Knightley, <that the same sort of handwriting often prevails in a family; and where the same master teaches, it is natural enough. But for that reason, I should imagine the likeness must be chiefly confined to the females, for boys have very little teaching after an early age, and scramble into any hand they can get. Isabella and Emma, I think, do write very much alike. I have not always known their writing apart.>

<Yes,> said his brother hesitatingly, <there is a likeness. I know what you mean—but Emma's hand is the strongest.>

<Isabella and Emma both write beautifully,> said Mr Woodhouse; <and always did. And so does poor Mrs Weston>—with half a sigh and half a smile at her.

<I never saw any gentleman's handwriting>—Emma began, looking also at Mrs Weston; but stopped, on perceiving that Mrs Weston was attending to some one else—and the pause gave her time to reflect, <Now, how am I going to introduce him?—Am I unequal to speaking his name at once before all these people? Is it necessary for me to use any roundabout phrase?—Your Yorkshire friend—your correspondent in Yorkshire;—that would be the way, I suppose, if I were very bad.—No, I can pronounce his name without the smallest distress. I certainly get better and better.—Now for it.>

Mrs Weston was disengaged and Emma began again—<Mr Frank Churchill writes one of the best gentleman's hands I ever saw.>

<I do not admire it,> said Mr Knightley. <It is too small—wants strength. It is like a woman's writing.>

This was not submitted to by either lady. They vindicated him against the base aspersion. <No, it by no means wanted strength—it was not a large hand, but very clear and certainly strong. Had not Mrs Weston any letter about her to produce?> No, she had heard from him very lately, but having answered the letter, had put it away.

<If we were in the other room,> said Emma, <if I had my writing-desk, I am sure I could produce a specimen. I have a note of his.—Do not you remember, Mrs Weston, employing him to write for you one day?>

<He chose to say he was employed>—

<Well, well, I have that note; and can shew it after dinner to convince Mr Knightley.>

<Oh! when a gallant young man, like Mr Frank Churchill,> said Mr Knightley dryly, <writes to a fair lady like Miss Woodhouse, he will, of course, put forth his best.>

Dinner was on table.—Mrs Elton, before she could be spoken to, was ready; and before Mr Woodhouse had reached her with his request to be allowed to hand her into the dining-parlour, was saying—

<Must I go first? I really am ashamed of always leading the way.>

Jane's solicitude about fetching her own letters had not escaped Emma. She had heard and seen it all; and felt some curiosity to know whether the wet walk of this morning had produced any. She suspected that it had; that it would not have been so resolutely encountered but in full expectation of hearing from some one very dear, and that it had not been in vain. She thought there was an air of greater happiness than usual—a glow both of complexion and spirits.

She could have made an inquiry or two, as to the expedition and the expense of the Irish mails;—it was at her tongue's end—but she abstained. She was quite determined not to utter a word that should hurt Jane Fairfax's feelings; and they followed the other ladies out of the room, arm in arm, with an appearance of good-will highly becoming to the beauty and grace of each.